\chapter{Introdução}

O controle patrimonial constitui uma atividade essencial em instituições públicas, pois está diretamente relacionado à rastreabilidade dos bens permanentes, à prestação de contas e ao uso adequado dos recursos institucionais. Em ambientes universitários, essa atividade tende a apresentar maior complexidade devido à quantidade de equipamentos distribuídos em laboratórios, salas administrativas e demais espaços de uso compartilhado. Nesse contexto, processos baseados em conferência manual ou em identificação visual individual podem gerar atrasos, inconsistências e dificuldades para manter o inventário físico alinhado aos registros administrativos \cite{tceba2018}.

No Colegiado de Ciência da Computação da Universidade Estadual de Santa Cruz, a gestão de ativos envolve equipamentos que podem ser utilizados por diferentes usuários e em diferentes ambientes. A ausência de mecanismos automatizados de identificação e atualização dificulta a verificação frequente da localização dos bens e torna o processo de inventário dependente de intervenção humana. Esse cenário é compatível com limitações observadas em estudos sobre gestão patrimonial e controle de estoques, nos quais métodos tradicionais, como conferência visual e código de barras, são apontados como menos eficientes em contextos com grande volume de itens \cite{brito2019}.

A Identificação por Radiofrequência (RFID) apresenta-se como alternativa para automatizar parte desse processo. Diferentemente de soluções ópticas, a tecnologia permite que etiquetas sejam reconhecidas por leitores compatíveis, sem a mesma dependência de contato visual direto. Essa característica favorece aplicações de inventário, rastreabilidade e auditoria, especialmente quando integrada a sistemas de informação capazes de armazenar, processar e apresentar os eventos de leitura \cite{alwadi2017,mashayekhy2022}.

Este trabalho propõe e implementa um protótipo funcional de software para inventário patrimonial baseado em RFID, voltado ao contexto do Colegiado de Ciência da Computação. A solução foi concebida para receber eventos gerados por diferentes dispositivos de captura por meio de uma API padronizada, preservando o fluxo previsto para sensores, gateways ou leitores em rede. A validação física concentrou-se no hardware disponível para o experimento, um leitor RFID de proximidade de baixo custo, utilizado para comprovar o funcionamento do núcleo computacional da proposta.

Essa delimitação evita atribuir ao protótipo resultados que dependeriam de uma infraestrutura RFID mais robusta. O foco está na verificação funcional do software e de sua capacidade de receber eventos de leitura, processá-los e refletir seus efeitos no inventário patrimonial. Essa distinção é coerente com a separação entre validação funcional de sistemas e avaliação completa de desempenho em ambiente operacional \cite{ieee1012}. Assim, a contribuição principal do trabalho está na arquitetura e na implementação do sistema, bem como na demonstração de sua viabilidade em um cenário controlado.

\section{Objetivos}

\subsection{Objetivo Geral}

Desenvolver um protótipo funcional de sistema web para inventário patrimonial baseado em RFID, capaz de registrar ativos, leitores, locais e eventos de leitura, permitindo a atualização do estado físico dos bens e o apoio à auditoria patrimonial no Colegiado de Ciência da Computação da UESC.

\subsection{Objetivos Específicos}

Para alcançar o objetivo geral e responder ao problema de pesquisa proposto, tornou-se necessário estabelecer etapas específicas de análise, desenvolvimento e validação. Nesse sentido, os objetivos específicos orientam o percurso metodológico do trabalho, ajudam a organizar as discussões e direcionam a investigação de forma mais precisa. Assim, foram definidos os seguintes objetivos específicos:

\begin{enumerate}
    \item Levantar os requisitos essenciais para o controle de ativos patrimoniais no contexto do colegiado.
    \item Projetar uma arquitetura modular para recebimento e processamento de eventos RFID.
    \item Implementar um protótipo web para registro, acompanhamento e auditoria de ativos patrimoniais.
    \item Integrar o protótipo a um leitor RFID de proximidade por meio de uma interface padronizada de eventos.
    \item Validar funcionalmente o fluxo principal de leitura, processamento e atualização do inventário.
\end{enumerate}

\section{Motivação e Justificativa}
\label{sec:justificativa}

A motivação deste trabalho decorre da necessidade de reduzir a dependência de processos manuais no controle patrimonial. A conferência individual de bens demanda tempo, atenção e disponibilidade operacional, além de estar sujeita a erros de registro, esquecimento de atualização e divergências entre a localização real dos itens e a localização registrada. Em ambientes acadêmicos, essas dificuldades são ampliadas pela circulação de equipamentos entre laboratórios, setores e responsáveis \cite{tceba2018,brito2019}.

A adoção de RFID no contexto de inventário patrimonial permite explorar uma forma de identificação automatizada e integrada a sistemas de informação. Estudos sobre RFID em gestão de estoques e patrimônio indicam ganhos potenciais de rastreabilidade, redução de falhas humanas e aumento de confiabilidade nas informações registradas \cite{brito2019,mashayekhy2022}. Além disso, arquiteturas baseadas em eventos, APIs e camadas intermediárias favorecem a integração entre dispositivos físicos e aplicações web, aproximando a solução de abordagens de Internet das Coisas \cite{abijaude2017,razzaque2015}.

Do ponto de vista institucional, o protótipo contribui como base tecnológica para modernizar o controle patrimonial do colegiado. Do ponto de vista acadêmico, o trabalho integra conceitos de desenvolvimento web, modelagem de dados, APIs REST, processamento de eventos e validação experimental com hardware acessível. A escolha por um leitor de proximidade de baixo custo delimita a validação física, mas mantém a arquitetura preparada para dispositivos RFID mais robustos, desde que estes comuniquem seus eventos à API definida. Essa cautela é relevante porque implantações RFID podem ser afetadas por fatores de ambiente, componentes e interferência, exigindo ensaios específicos para avaliação física da solução \cite{knapp2022}.

Quanto à organização do texto, o Capítulo 2 apresenta a revisão de literatura sobre controle patrimonial, RFID, Internet das Coisas, auditoria e trabalhos relacionados. O Capítulo 3 descreve os materiais, os requisitos, a modelagem, a arquitetura e os procedimentos metodológicos adotados. O Capítulo 4 apresenta a implementação do protótipo e os resultados obtidos na validação funcional. Por fim, o Capítulo 5 reúne as conclusões, as limitações observadas e as possibilidades de continuidade do trabalho.
